% Continuation of the grouped manuscript on Boolean Doob chains.
\section{Controlling the denominator by the weight enumerator}
\label{sec:enumerator}
Let \(C=\operatorname{im}(A^\top)\subseteq\mathbb F_2^L\) have rank \(r\).
Assume \(E\pi=1\), \(\|\pi\|_2\le B\), and marginal matching through order
\(k\to\infty\). The density is strictly positive in each dimension.
Define the weighted counts
\[
 Z=\sum_{w\in C}2^{-|w|},\qquad T_k=\sum_{w\in C,\ |w|>k}2^{-|w|}.
\]
The first counts all available characters; the second uses only those
whose coefficients may survive matching.

\paragraph{Two criteria.}
If \(2^{-k} Z\to0\) and \(T_k\to0\), then the entire ordered spectrum
recovers, for any linear quotient. If the columns instead contain two
disjoint bases, \(2^{-k} Z\to0\) alone suffices. In this second case
exponential decay of \(2^{-k} Z\) gives exponential recovery in \(k\).
In particular a sufficient rank condition for exponential recovery is
\[
 \limsup r/k<1/\log_2(5/4).
 \tag{W1}
\]
There is no bound on the peak or inverse of the stationary density among
these hypotheses. The constants from the ordinary noise operator act on
the full cube, including eigenvectors that vary within a quotient fiber.

\paragraph{Pointwise denominator control.}
Put \(F=Qh,\ a=Eh,\ V=\|F-a\|_2^2=1-a^2\).
Cauchy--Schwarz gives \(1=(Ehg)^2\le a E(\pi g)\).
The nonconstant coefficients of \(\pi\) vanish through degree \(k\), and
\(g=QF\). By Parseval,
\(V\le a\|\pi-1\|_2\,2^{-(k+1)/2}\sqrt V\).
This includes \(V=0\). Since \(F\) is a quotient function, Cauchy--Schwarz
in its Walsh expansion gives
\[
 \|g-1\|_\infty\le V+\sqrt{Z-1}\sqrt V
 \le B^2 2^{-(k+1)} + B\,2^{-(k+1)/2}\sqrt{Z-1}=:d_k .
 \tag{W2}
\]
Here \(1-a\le V\) since \(0<a\le1\).
In particular \(g\ge 1-d_k\) on every vertex. Any exceptional set defined
by a fixed denominator cutoff below one is eventually empty when
\(2^{-k} Z\to0\).

\paragraph{Lifting the signed comparison.}
Write \(H=\|QM_{\pi-1}Q\|\) on the full cube. For a character \(\chi_w\),
\(QM_{\chi_w}Q\) is a weighted permutation in the orthonormal Walsh basis,
with largest entry \(2^{-|w|/2}\): its entries have exponents
\(-(|u|+|u+w|)/2\) and \(u=0\) attains the maximum. Thus
\(H\le\|\pi-1\|_2\sqrt{T_k}\). When \(d_k\le1/2\), positivity of \(\pi\)
and (W2) give a quadratic-form bound
\[
 \|QM_{U-\pi}Q\|\le 8d_k\,\|QM_\pi Q\|\le8d_k(1+H).
\]
Indeed \(|g^{-2}-1|\le 8d_k\) in that range.
The two Gram products from the preceding section and min--max now give
\(\Delta\le H+8d_k(1+H)\). This proves the first criterion.

For two bases there is a second bound on \(H\) irrespective of \(T_k\).
Take \(\lambda=3/4,\ \theta=1/(\sqrt2\lambda)<1\), so
\(Q=P_\theta P_\lambda\). On the quotient the two full bases give
\(P_\lambda:L^2\to L^4\) since \(1+\lambda^{-4}>4\).
The positive leading eigenvector argument for
\(P_\lambda M_{|\pi-1|}P_\lambda\) lifts its norm bound
\(\|\pi-1\|_2\) to the full cube and hence bounds the signed form too.
Projecting at degree \(j=\lfloor k/2\rfloor\) on both sides, the
low--low block of \(QM_{\pi-1}Q\) vanishes by matching; each other block
has at least one \(P_\theta\) factor restricted to degrees greater than
\(j\). Therefore \(H\le3B\theta^{j+1}\). This proves the second criterion.

For disjoint bases \(J_1,J_2\), \(|w|\ge |w_{J_1}|+|w_{J_2}|\).
Each restriction is a bijection from \(C\) to the binary \(r\)-cube.
Cauchy--Schwarz in the unnormalized sum gives
\(Z\le(1+2^{-2})^r\), proving (W1). More generally a fractional packing
of weight \(t>0\) gives \(Z\le(1+2^{-t})^r\): apply generalized
H\"older with weights \(b_J/t\) to the factors \(2^{-t|w_J|}\) for the
positive base weights \(b_J\), using the column capacity constraint.
This last counting bound alone does not supply the two-base signed estimate.

\paragraph{An example beyond peak control.}
For comparison, a direct truncation gives a sufficient peak rate
\(\limsup k^{-1}\log_2\|\pi\|_\infty<7/15\) for two bases at \(p=2\).
Indeed set \(s=5/3\); \(Q:L^2\to L^5\) on the quotient gives the
positive multiplier bound from \(L^s\). The preceding section bounds
\(\|g^{-1}\|_q\) for every \(q<15\).
For \(E=\{|g-1|>2^{-\beta k}\}\), its matching estimate gives
\(\mu(E)\le C_B 2^{-(1-2\beta)k}\).
Then \(\|U{\bf1}_E\|_s\le \|\pi\|_\infty\|g^{-1}\|_q^2
\mu(E)^{3/5-2/q}\). Taking \(q\) close to 15 and \(\beta>0\) small
proves exponential decay. The \(\pi{\bf1}_E\) term costs
\(B\mu(E)^{3/5-1/2}\), and off \(E\) the relative error is
\(O(2^{-\beta k})\). The signed comparison is as above.

Take odd \(k\ge3\) sufficiently large, \(r=3k\), and two identical
copies of each of the \(r\) coordinate columns. A \(k\)-dimensional
subspace \(D\subseteq\mathbb F_2^r\) exists with minimum weight precisely
\((k+1)/2\). Fix the first generator to have that weight, then choose
\(k-1\) independent uniformly sampled generator vectors. Every combination
using a sampled vector is uniform, so the probability of one of these
combinations having weight \(\le(k-1)/2\), including zero, is at most
\[
 2^{k-r}\sum_{i\le(k-1)/2}\binom ri
 \le 2^{-k[2-3H_2(1/6)]}<1 .
\]
The bound includes linear dependence. Here \(H_2\) is binary entropy and
\(6^6<16\cdot5^5\) verifies the last strict inequality. Hence such a choice
exists for each \(k\).

Let \(S=D^\perp,\ u=2^{-k}\) and
\(\pi_0=1+u^{-1/2}({\bf1}_S-u)\) on the quotient.
To ensure intrinsic rank \(r\) and permit a vanishing minimum, choose
\(y_0\notin S\), put \(z=(1,\ldots,1)\), \(\epsilon=2^{-2k}\), and set
\[
 W(y)=\frac{\chi_z(y-y_0)+\sum_{i=1}^{r}\chi_{z+e_i}(y-y_0)}{r+1},
 \qquad
 \pi=\pi_0-(1-\sqrt u-\epsilon)W .
\]
The mean is one and \(\min\pi=\epsilon\) since \(|W|\le1\) and
\(W(y_0)=1\). Its second moment stays bounded; on mass \(u\) the height
is comparable to \(u^{-1/2}\), so every fixed moment above two diverges
and \(k^{-1}\log_2\|\pi\|_\infty\to1/2\). All characters in \(D\)
have ambient weight twice their quotient weight. The fixed shortest
word survives because the added characters have quotient weights \(r,r-1\).
The largest matching order is therefore exactly \(k\).
All \(r+1\) added characters have nonzero coefficients even if they belong
to \(D\): the original modulus there is \(\sqrt u\), exponentially smaller
than the added modulus. Their differences span the quotient. Its intrinsic
packing is exactly two by column count. Positive symmetric scaling gives
the corresponding Doob chain in every finite dimension.
Here \(Z=(5/4)^{3k}\) exactly and \(2^{-k} Z=(125/128)^k\).
Thus (W2) controls every vertex of a family that exceeds the peak rate
at its actual matching order and has no bounded higher stationary moment.

\paragraph{Scope.}
Weight enumerators and noisy functions on binary linear codes have a
substantial prior literature, including
\href{https://arxiv.org/abs/1809.09696}{Samorodnitsky's noisy-function
inequality} and his \href{https://arxiv.org/abs/2205.02051}{work on random
linear codes}. The random-code volume estimate and Walsh multiplier
identities above are elementary inputs. The conclusion here couples them
to stationary scaling and the ordered ambient spectrum, via (W2).
The rank ceiling is sufficient. This calculation leaves open whether two
bases and a bounded second moment suffice at unrestricted rank.
